\figrel{fig01.png}{E04. Matrizes das duas configurações com extrator ImageNet congelado.}
\figrel{fig02.png}{E04. Perda ponderada no treino da cabeça e resumo da validação.}
\figrel{fig03.png}{E05. Ajuste seletivo. A escala linear agrupa perdas muito pequenas perto da origem.}
\figrel{fig04.png}{E06. Ablação das fases. A escala de 0 a 1 comprime diferenças próximas ao teto.}
\figrel{fig05.png}{E07.1. ImageNet versus três inicializações aleatórias com extrator fixo.}
\figrel{fig06.png}{E07.2. Treino integral do zero e referência E04.}
\figrel{fig07.png}{E07.3. CNN 1D e referência ResNet-18 com STFT do E07.2.}
\figrel{fig08.png}{E08.1. Estatísticas, CNN e fusão híbrida. Contagens e logloss.}
\figrel{fig09.png}{E09.1. Robustez do treino limpo e misto sob AWGN a 25 kHz.}
\figrel{fig10.png}{E09.2. Proporções de AWGN com orçamento constante.}
\figrel{fig11.png}{E09.3. Ganhos independentes e teste cruzado com AWGN.}
\figrel{fig12.png}{E09.4. Auditoria de suporte por variante do Gate 4.}
\figrel{fig13.png}{E09.5. Treino canônico versus deformação temporal.}
\figrel{fig14.png}{E09.6. Modulação e combinação comparadas com AWGN puro.}
\figrel{fig15.png}{E09.7. Adição de variantes versus repetição limpa pareada.}
\figrel{fig16.png}{E10.1. Realizações inéditas de AWGN a 25 kHz.}
\figrel{fig17.png}{E10.2. Pós-retenção simulado. Eixo de SNR nominal a 200 kHz.}
\figrel{fig18.png}{E10.3. Pós-retenção em arquivos do Gate 4.}
\figrel{fig19.png}{E10.4. Modelos ABC forçados sem B e alternativa AC.}
\figrel{fig20.png}{E10.5. Treino AWGN para a alternativa AC.}
\figrel{fig21.png}{E10.6. Varredura exploratória de qualidade para AC.}
\figrel{fig22.png}{E10.7. Bateria final de robustez e referências clássicas.}
\figrel{fig23.png}{E11. Oclusões de fases, tempo e frequência.}
\figrel{fig24.png}{E11. PCA ajustada no treino e projetada na validação.}
\figrel{fig25.png}{E12. Trajetória de escores em R03 sob condição limpa.}
\figrel{fig26.png}{E12. Trajetória de escores em R09 sob condição limpa.}
\figrel{fig27.png}{E12. Custo por janela, episódios de falso alarme e atraso com persistência de três decisões.}
\clearpage\figrel{mapas1.png}{E11. Grad-CAM e saliência. Parte 1.}
\clearpage\figrel{mapas2.png}{E11. Grad-CAM e saliência. Parte 2.}
\figrel{exemplo_pre.png}{E03. R03, janela pré-evento, dois ciclos. Painéis A/B/C da esquerda para a direita, com barra fixa à direita.}
\figrel{exemplo_pos.png}{E03. R03, janela pós-evento, dois ciclos. Painéis A/B/C da esquerda para a direita, com barra fixa à direita.}
